% GENERATED FILE. Edit canonical lessons, then run npm run generate:books.
% canonical-source-hash: fnv1a64:43507e1f31aef569
% canonical-lessons: TA-C526-tali, TA-C526-vatakku-saute, TA-C526-kilaru, TA-C526-turuvu, TA-C526-tuvu

\chapter{Cooking Steps}
\label{ch:ta-cooking-steps}
\hlchaptermodality{526}

\begin{chapteropening}
\textbf{By the end of this chapter:} \emph{I can describe how a South Indian dish is made.}

Complete the last lesson of chapter 526: I can describe how a South Indian dish is made.
\end{chapteropening}

\section[\texorpdfstring{\ta{தாளி} (tāḷi)}{தாளி (tāḷi)}]{\ta{தாளி} (tāḷi) — to temper spices}

\label{lesson:TA-C526-tali}



\begin{quote}
Before the new one: say the Tamil for to steam, then the Tamil for to be cooked.
\end{quote}

\subsection*{You'll want to know: \ta{தாளி}}
\textbf{\ta{தாளி}} — \emph{tāḷi} — \textquotedblleft{}to temper spices\textquotedblright{}.

\begin{quote}
\ta{அம்மா} \ta{கடுகையும்} \ta{கறிவேப்பிலையையும்} \ta{எண்ணெயில்} \ta{தாளித்தாள்}.
\end{quote}

\begin{quote}
\emph{ammā kaṭukaiyum kaṟivēppilaiyaiyum eṇṇeyil tāḷittāḷ.}
\end{quote}

\begin{quote}
Mother tempered mustard seeds and curry leaves in oil.
\end{quote}

Past (I …): \ta{தாளித்தேன்}, \emph{tāḷittēṉ}.

Spices fried in hot oil and poured over a dish. The noun is \ta{தாளிப்பு}.

\begin{cousinweb}
\textbf{\ta{தாளி}} (\emph{tāḷi}) is native Dravidian: DEDR 3186 gives \ta{தாளி} (\emph{tāḷi}), \textquotedblleft{}to season and flavour a curry with spices fried in ghee or oil\textquotedblright{}.
\end{cousinweb}

\begin{grammarlens}[title={What you've built}]
One more word for describing what you lived and felt.
\end{grammarlens}

\subsection*{Your turn}
\begin{itemize}
\raggedright
  \item \emph{Say it:} \emph{tāḷi}
  \item \emph{Say it:} \emph{tāḷi}, once more
  \item \emph{Say it:} \emph{ammā kaṭukaiyum kaṟivēppilaiyaiyum eṇṇeyil tāḷittāḷ.} — Mother tempered mustard seeds and curry leaves in oil.
  \item \emph{Recall:} say the Tamil for to steam, then the Tamil for to be cooked, then say \emph{tāḷi} again
\end{itemize}

\begin{tcolorbox}[breakable,colback=teal!4,colframe=teal!35!black,title={Before you move on}]
What does \ta{தாளி} mean? (\textquotedblleft{}To temper spices\textquotedblright{}.) Say it once more.
\end{tcolorbox}
\section[\texorpdfstring{\ta{வதக்கு} (vatakku)}{வதக்கு (vatakku)}]{\ta{வதக்கு} (vatakku) — to sauté; to fry lightly}

\label{lesson:TA-C526-vatakku-saute}



\begin{quote}
Before the new one: say the Tamil for to be cooked, then the Tamil for to temper spices.
\end{quote}

\subsection*{You'll want to know: \ta{வதக்கு}}
\textbf{\ta{வதக்கு}} — \emph{vatakku} — \textquotedblleft{}to sauté; to fry lightly\textquotedblright{}.

\begin{quote}
\ta{அவள்} \ta{வெங்காயத்தைப்} \ta{பொன்னிறமாக} \ta{வதக்கினாள்}.
\end{quote}

\begin{quote}
\emph{avaḷ veṅkāyattaip poṉṉiṟamāka vatakkiṉāḷ.}
\end{quote}

\begin{quote}
She sautéed the onions until golden.
\end{quote}

Past (I …): \ta{வதக்கினேன்}, \emph{vatakkiṉēṉ}.

Lighter than \ta{வறு} (\emph{vaṟu}), to roast, and \ta{பொரி} (\emph{pori}), to deep-fry.

\begin{cousinweb}
\textbf{\ta{வதக்கு}} (\emph{vatakku}) is native Dravidian: DEDR 5242 gives \ta{வதக்கு} (\emph{vatakku}), \textquotedblleft{}to roast, to dry\textquotedblright{}, beside \ta{வதங்கு}, \textquotedblleft{}to wither\textquotedblright{}.
\end{cousinweb}

\begin{grammarlens}[title={What you've built}]
One more word for describing what you lived and felt.
\end{grammarlens}

\subsection*{Your turn}
\begin{itemize}
\raggedright
  \item \emph{Say it:} \emph{vatakku}
  \item \emph{Say it:} \emph{vatakku}, once more
  \item \emph{Say it:} \emph{avaḷ veṅkāyattaip poṉṉiṟamāka vatakkiṉāḷ.} — She sautéed the onions until golden.
  \item \emph{Recall:} say the Tamil for to be cooked, then the Tamil for to temper spices, then say \emph{vatakku} again
\end{itemize}

\begin{tcolorbox}[breakable,colback=teal!4,colframe=teal!35!black,title={Before you move on}]
What does \ta{வதக்கு} mean? (\textquotedblleft{}To sauté; to fry lightly\textquotedblright{}.) Say it once more.
\end{tcolorbox}
\section[\texorpdfstring{\ta{கிளறு} (kiḷaṟu)}{கிளறு (kiḷaṟu)}]{\ta{கிளறு} (kiḷaṟu) — to stir a pot; to stir up}

\label{lesson:TA-C526-kilaru}



\begin{quote}
Before the new one: say the Tamil for to temper spices, then the Tamil for to sauté.
\end{quote}

\subsection*{You'll want to know: \ta{கிளறு}}
\textbf{\ta{கிளறு}} — \emph{kiḷaṟu} — \textquotedblleft{}to stir a pot; to stir up\textquotedblright{}.

\begin{quote}
\ta{பாயசம்} \ta{அடிப்பிடிக்காமல்} \ta{இருக்க} \ta{அம்மா} \ta{அதைக்} \ta{கிளறிக்கொண்டே} \ta{இருந்தாள்}.
\end{quote}

\begin{quote}
\emph{pāyacam aṭippiṭikkāmal irukka ammā ataik kiḷaṟikkoṇṭē iruntāḷ.}
\end{quote}

\begin{quote}
Mother kept stirring the payasam so that it would not stick.
\end{quote}

Past (I …): \ta{கிளறினேன்}, \emph{kiḷaṟiṉēṉ}.

Close to \ta{கலக்கு} (\emph{kalakku}), to mix; \ta{கிளறு} is the stirring of something thick with a ladle.

\begin{cousinweb}
\textbf{\ta{கிளறு}} (\emph{kiḷaṟu}) is native Dravidian: DEDR 1588 gives \ta{கிளறு} (\emph{kiḷaṟu}), \textquotedblleft{}to poke, to dig up, to stir with a ladle\textquotedblright{}.
\end{cousinweb}

\begin{grammarlens}[title={What you've built}]
One more word for describing what you lived and felt.
\end{grammarlens}

\subsection*{Your turn}
\begin{itemize}
\raggedright
  \item \emph{Say it:} \emph{kiḷaṟu}
  \item \emph{Say it:} \emph{kiḷaṟu}, once more
  \item \emph{Say it:} \emph{pāyacam aṭippiṭikkāmal irukka ammā ataik kiḷaṟikkoṇṭē iruntāḷ.} — Mother kept stirring the payasam so that it would not stick.
  \item \emph{Recall:} say the Tamil for to temper spices, then the Tamil for to sauté, then say \emph{kiḷaṟu} again
\end{itemize}

\begin{tcolorbox}[breakable,colback=teal!4,colframe=teal!35!black,title={Before you move on}]
What does \ta{கிளறு} mean? (\textquotedblleft{}To stir a pot; to stir up\textquotedblright{}.) Say it once more.
\end{tcolorbox}
\section[\texorpdfstring{\ta{துருவு} (turuvu)}{துருவு (turuvu)}]{\ta{துருவு} (turuvu) — to grate; to scrape}

\label{lesson:TA-C526-turuvu}



\begin{quote}
Before the new one: say the Tamil for to sauté, then the Tamil for to stir a pot.
\end{quote}

\subsection*{You'll want to know: \ta{துருவு}}
\textbf{\ta{துருவு}} — \emph{turuvu} — \textquotedblleft{}to grate; to scrape\textquotedblright{}.

\begin{quote}
\ta{பாட்டி} \ta{தேங்காயைத்} \ta{துருவினாள்}.
\end{quote}

\begin{quote}
\emph{pāṭṭi tēṅkāyait turuviṉāḷ.}
\end{quote}

\begin{quote}
Grandmother grated the coconut.
\end{quote}

Past (I …): \ta{துருவினேன்}, \emph{turuviṉēṉ}.

Grated coconut is \ta{தேங்காய்த்} \ta{துருவல்}.

\begin{cousinweb}
\textbf{\ta{துருவு}} (\emph{turuvu}) is native Dravidian: DEDR 3339 gives \ta{துருவு} (\emph{turuvu}), \textquotedblleft{}to bore, to scrape out the pulp of a coconut\textquotedblright{}.
\end{cousinweb}

\begin{grammarlens}[title={What you've built}]
One more word for describing what you lived and felt.
\end{grammarlens}

\subsection*{Your turn}
\begin{itemize}
\raggedright
  \item \emph{Say it:} \emph{turuvu}
  \item \emph{Say it:} \emph{turuvu}, once more
  \item \emph{Say it:} \emph{pāṭṭi tēṅkāyait turuviṉāḷ.} — Grandmother grated the coconut.
  \item \emph{Recall:} say the Tamil for to sauté, then the Tamil for to stir a pot, then say \emph{turuvu} again
\end{itemize}

\begin{tcolorbox}[breakable,colback=teal!4,colframe=teal!35!black,title={Before you move on}]
What does \ta{துருவு} mean? (\textquotedblleft{}To grate; to scrape\textquotedblright{}.) Say it once more.
\end{tcolorbox}
\section[\texorpdfstring{\ta{தூவு} (tūvu)}{தூவு (tūvu)}]{\ta{தூவு} (tūvu) — to sprinkle; to scatter}

\label{lesson:TA-C526-tuvu}



\begin{quote}
Before the new one: say the Tamil for to stir a pot, then the Tamil for to grate.
\end{quote}

\subsection*{You'll want to know: \ta{தூவு}}
\textbf{\ta{தூவு}} — \emph{tūvu} — \textquotedblleft{}to sprinkle; to scatter\textquotedblright{}.

\begin{quote}
\ta{சமையல்காரர்} \ta{கொத்தமல்லியை} \ta{மேலே} \ta{தூவினார்}.
\end{quote}

\begin{quote}
\emph{camaiyalkārar kottamalliyai mēlē tūviṉār.}
\end{quote}

\begin{quote}
The cook sprinkled coriander on top.
\end{quote}

Past (I …): \ta{தூவினேன்}, \emph{tūviṉēṉ}.

What is sprinkled takes -\ta{ஐ}: \ta{உப்பைத்} \ta{தூவு}, to sprinkle salt.

\begin{cousinweb}
\textbf{\ta{தூவு}} (\emph{tūvu}) is native Dravidian: DEDR 3394 gives \ta{தூவு} (\emph{tūvu}), \textquotedblleft{}to sprinkle, to strew, to scatter\textquotedblright{}.
\end{cousinweb}

\begin{grammarlens}[title={What you've built}]
One more word for describing what you lived and felt.
\end{grammarlens}

\subsection*{Your turn}
\begin{itemize}
\raggedright
  \item \emph{Say it:} \emph{tūvu}
  \item \emph{Say it:} \emph{tūvu}, once more
  \item \emph{Say it:} \emph{camaiyalkārar kottamalliyai mēlē tūviṉār.} — The cook sprinkled coriander on top.
  \item \emph{Recall:} say the Tamil for to stir a pot, then the Tamil for to grate, then say \emph{tūvu} again
\end{itemize}

\begin{tcolorbox}[breakable,colback=teal!4,colframe=teal!35!black,title={Before you move on}]
What does \ta{தூவு} mean? (\textquotedblleft{}To sprinkle; to scatter\textquotedblright{}.) Say it once more.
\end{tcolorbox}
